% $Id: finalreport.tex 4414 2013-02-11 10:48:35Z jabriffa $

\documentclass{surreydissertation}

% LOAD MISSING PACKAGES HERE
\usepackage{amsmath}  % Fixes the 'cases' and '\text' errors
\usepackage{amsfonts} % Provides the blackboard bold \mathbb font
\usepackage{hyperref} % Put this in the preamble

% Write the approved title of your dissertation
\title{Brain-Inspired Machine Unlearning: Emulating the Transformative Success of Emotional Memory Reconsolidation Therapy}

% Write your full name, as in University records
\author{Ansen Lam}

% Write your URN, as in University records
\urn{6842754}

% Write the month and year of your submission
\date{Aug 10, 2026}

% Uncomment the line for the Degree you are registered for
%\degree{Bachelor of Science in Computer Science}
\degree{Bachelor of Science in Computing and Information Technology}

% Write the full name of your supervisor, without any titles
\supervisor{Dr Asgha, Rizwan}

\begin{document}
\maketitle
\tableofcontents

\chapter{Abstract}
Reconsolidation is a replicable process in neurobiology to treat memory recall that is emotionally disruptive. It works on the basis of recreating the context of the surrounding trigger for the emotion, and introducing a new safer emotion to overwrite this trigger. This effect is long lasting and permanent with proper treatment. This suggests that memories, which is widely agreed to form in the hippocampal formation, indexed through time in it's growth, and gradually migrates to the outer cortex, is not permanent.

This suggests a worthy subject of intrigue in which an engantled neuro network model exibits the ability to cleanly "forget" or "rewrite" influences from past examples, as if an indexed storage.

Inspired by a distributive network able to cleanly target a memory and modify it, we investigated the current approaches of unlearning algorithms, Hopfield network associative memory storage, abstract models of latent cause theory that describe reconsolidation therapy, processes of hippocampal indexing like pattern completion and seperation, and adapted the 2023 unlearning competition metric to measure an attempt at a new classification unlearning algorithm with these components and compared against other approaches.

\chapter*{Acknowledgements}
    Dr Asgha, Rizwan, for his guidance and tolerance.
\chapter{Introduction}
Machine unlearning concerns the removal of data trained into a model, with rich regulatory and law implications. The most notable refers to the UK General Data Protection Regulation Article 17 (the right to forget), protecting data ownership, consent, personal safety among others~\cite{edps2025,liu2024}. This task is incredibly daunting due to the entangled sequential transformation during training, akin to a one-way cryptographic function~\cite{cao2015}. To date, computationally acceptable solutions are mere approximations with degrees of residual influence, while exact solutions are far from feasibly operational~\cite{xu2023,wang2024}.

Memory reconsolidation research might provide a significant piece to the puzzle with their profound breakthrough on emotional memory therapy, breaking the long standing understanding that emotional memory consolidation is permanent through synaptic lock~\cite{brunet2008,brunet2014,brunet2018,wright2021}. Most significant is how clean and permanent this form of therapy effected memory~\cite{kamboj2020}.

Emulating this in context of a learnable model would have profound implications. Firstly, the recall and memory requirement for re-consolidation would suggest that a data address relationship exists in such models~\cite{ramsauer2020}. Secondly, reconsolidation allows for memory alteration, allowing data dynamic models. And thirdly, this data-targeted approach is likely computationally feasible.

The following sections discusses the main ideas behind the inspired unlearning algorithm.

\section{Machine unlearning models and techniques}
Machine unlearning techniques have been developed to address a range of distinct technical problems. A systematic compilation and analysis of these problems and their corresponding solutions provides a useful basis for examining potential analogies with the brain’s capacity for selective forgetting.

\subsection{Exact unlearning}  
These methods aim to remove the influence of the forget set as if the model has been completely retrained. They limit data influence by partitioning the dataset into shards, trained into sub-models, forming an ensemble. Unlearning is achieved by retraining,only the affected sub-models, reducing computational cost. Although exact, these methods are still impractical in scaling, accuracy degredation and design complexity. SISA \cite{bourtoule2021sisa}, ARCANE \cite{arcane2022} and others follow this methodology

Since only the data affected partition is retrained, the forget set influence is physically bounded from other possible entanglement, this makes influences contained. Any model that is not exact is approximate.

However most unlearning problems doesn't have known partition of dataset, often at the request after the fact of training. It also assumes the often wrong assumption that the final model does not have dependent data relavant to the output between partitioned models. Finally, models that think alike often come to the same conclusion, so however exact this unlearning method is, it cannot guarantee that outputs resemble that of the request take down based on convergence evolution.

\subsection{Direct Fine-tuning}
Direct fine-tuning methods perform additional training steps on the forget set, like fine-tuning. These methods are computationally simple, but is prone to cause unintentional damage existing knowledge,akin to catastrophic forgetting in continual learning. Below are some examples:

Gradient ascent (GA), additionally train a loss over the forget set using negative log-likelihood loss, To maintain model accuracy, regularization loss over the retrain set is often used, like a gradient descent or KL divergence. 

Reinforcement learning (RL) methods like Proximal Policy Optimization (PPO), Preference optimization (PO), Direct Preference optimization (DPO), negative preference optimization (NPO) that treats the forget set as preferences to policy on the model.

These models are fast approximates, however they do not solve the fundamental entangled nature of trained data by inherently using similar methods for training to untrain it's data, introducing the same residual influence or catastrophic forgetfulness with the inherent setup.

\subsection{Localized Parameter Modification}
Localized Parameter Modification targets and adjusts specific parts of the model responsible for the forget set. These methods are precise, but require a lot of assumptions on how knowledge is located in the model. For example, language models represent tokens as vectors, and training tokens to point to a space of misunderstanding on the forget set of harmful content encodes censorship.

These methods solve direct fine-tunings' referenceless re-training by defining one in output, however model entanglment persists, and assumes all possible unlearning concepts/meanings as the defined structure without higher order combinations, or emergence behaviour of dependent meanings.

\subsection{Input/Output-based unlearning}
Input/Output-based unlearning focuses filtering the model's inputs or outputs without changing the core model to achieve unlearning on the forget set. 

This is the next-best strategy used in most commercial AI outputs due to it's simplicity and speed, however due to it's ad hoc nature, it introduces manual error, one size fits all censorship, and addes maintanance and ill-automated scaling for highly related (or not) cases.

\section{Metric on machine unlearning}
A good machine unlearning algorithm should be power efficient, behave as if without prior knowledge of the unlearnt, and retain it's model performance.

The 2023 Google unlearning competion benchmark, provides a starting point to measure classification models. It assesses the performance of approximated unlearning algorithms on classification models. It has three metrics: forgetting quality, efficiency and utility. The bases for comparison is inspired by Differential Privacy, which is a measure on what an adversary can learn about the leak of a data point. It measures the change of distribution regarding other datapoints if one is leaked.

To reframe the quality (unlearning) measure to unlearning, an advasary, in this case the auditor model, measures unlearning by comparing the distributional change of the unlearned model against the base completely retrained model.

It trains a number of baseline models (full re-trained without the forget set), and a number of unlearning models. By giving these modles a random seed, they exhibit a statistical model for each input. By performing a membership inference attack, where an example was given a probablity of being in the training data, these statistical models are compared for similarity, giving a score to quality (unlearning).

This metric itself is not time efficient, the number of models that needs to be trained to have a meaningful statistical comparison of it's variations is far higher than the infesible scaling of retraining from scratch. It does not make sense to create an efficient unlearning algorithm beyond exact retraining, only to need to verify it with ten times or more compute on that baseline. However for small models this provides a useful metric prior to scaling.


\section{Memory consolidation and emotional memory therapy}
Memory consolidation is the process of gradual migration and stabilization of memory from the hippocampus to the outer cortex. It's understanding mainly comes from protein inhibitor experiments, where the hippocampal formation is found to have a central role in all forms of memories. This process is long thought to be permanent (synaptic lock theory), until it was disproven by a shock (unconditional stimulus US) and queue (conditional stimulus CS) experiment on rats. Protein inhibitors administered after the event of recalling a memory is shown to disrupt the memory permanently, this event is called reconsolidation, in which the memory is liable for update.

While reconsolidation is successfully implemented in treating emotional memory, it's replicability comes under question in other forms of memory, leading to several proposals to fix it. One theory of interest is latent cause theory \cite{gershman2017computational}, which incoporates observations from behavioral experiments, like conditioning, extinction, and renewal paradigms. It proposes that the brain learns the probability of an event generated by a latent cause, and memory is thus the association between them. It serves as a unifying theory for a wide range or observations following reconsolidation.

It models a three way model of three parts: CS given cause, US given cause, and cause given prior causes.
CS given cause is drawn from the product normal distribution of it's expected values of that cause.
US given cause id drawn from the normal distribution in the mean of the hebbian sum of weights and CS.
cause given prior causes is drawn from a temporal kernal, which implements the continguity principle of temporal proximity (similar events in similar time frame) infers cause proximity.

This model assumes that a cause is always present in any stimuli, and operates on a very abstract level. Recursively, causes can be theoretically represented by epsodic neurons which are neurons found to trigger only upon a singular experience or concept, linking through them with their pattern expressions in the outer cortex upon firing. This is potentially apt in explaining how and what conditions do these "pointer" neurons form, and is a good step to figure out the next step of how these "pointers" gradually migrate to the outer cortex for long term memory or reflex.

For unlearning, these pointers can be formed in a multitude of layers, with specific data, experience, and concepts. Formed through a common pointer, they could be altered to give us a targeted data-tied expression of a model which takes inputs from these "pointers". This would in turn greatly simplify the problem of engantled influence by being a representative of unwanted patterns, or indeed any pattern of interest, giving meaning to parameters in a model.

\section{Hopfield network, Associative memory modeling}
Hopfield networks are originally inspired by the numerous stable states in a ferromagnet's spin configurations, or Ising spins. This model is adapted to fully interconnected neurons where an update function changes a neurons' state based on the weighted sum of the others. Unlike typical forward networks, this network is fully recurrent and it's memory reliant on the convergence of the update function. These stable points are analogous to memory patterns. The encoded patterns, or the specific stable points, are dictated by an interaction function to randomly update neurons based on the weights connected to it. Numerous improvements to It has since been made to it including:
exponential pattern storage capacity relative to number of neurons

maintaining exponential capacity with standard two-body synnaptic connections by adding one hidden layer (one input and one output), giving the same capacity in realistic neuron configurations.(cite)

Generalizing interaction function by making it differentiable in a propagation layer, with an update function equivalent to transformer attention heads. The transformer equivalence states that if learned Q, K, V projects the query (input pattern), key (the current pattern), and value (the result pattern after going through update), the modern Hopfield network is equivalent to an attention head in a transformer~\cite{ramsauer2020}.

Hopfield networks is a literal translation of neural configurations in the brain to a computational model.It is often used in modeling associative memory due to their partial cue pattern completion property, where patterns in small sections could be rapidly auto-completed to the full pattern by approximation.


\section{Hippocampal index theory}
The theory states that the hippocampal formation stores pointers to access memories in the outer cortex, as such the full context of the memory are integrated as one. These memories gradually stabilize in the cortex. Two proceses underly it's function: pattern seperation and pattern completion. Pattern completion concerns the brain's ability to imagine the full context of a memory, or a predictoin of a memory, by being exposed to a partial cue. This process is mainly attributed to CA3, a section of the hippocampus.  Pattern seperation is where similar experiences are orthogonalized into different memories. Originating from EC,  it is routed to DS (Dentate Gyrus) for sparse encoding, then it is stored in CA3, then routed through CA to the cortical memories.

These two processes often complement each other in the Hippocampus, where pattern completion is mainly attributed to DG-CA3 associations, and pattern completion CA3-CA1 and CA1 to subiculum connections. This is supported by activity observations, in anterior hippocampus, where it is more active in the initial stages of information extraction; and activity in posterior hippocampus, which is more active in maintaining representations.

Drawing on the interpretation that these engram cell encods latent cause on signal expression in the outer cortex, the hippocampal index theory suggests that we could populate the causes of expressions with pattern seperation, thereby controlling how fine-grained our perceptive cause are. At the same time suggest that we can dictate relevant causes in a short period of time using pattern completion, to filter out most likely causes. This is relevant in time and energy savings that greatly reduces the amount of causes considered.

% \section{key problems and challenges}

% \section{Biological accuracy vs property emulation}
% When bridging the gap between structure of atomic scale to properties on the neural system level, biological accuracy, is often sacrificed for computaionala efficiency.

% \section{opposite problem nature of brain aligned memory systems}
% When studying systems regarding memory models of the human brain, the problem turns into that of, maintaining the ability to obtain new information, while keeping the previous acquired memory intact. This is inherently opposite of unlearning. where utility is kept on the basis that the information in question is removed.

% \subsection{past work and state-of-the-art}
% (HICL achieved FLOPS advantage and performance advantage over other continual learning models)
% what other work has been done on it or aspects relating to,
% key technologies and state-of-the-art

\chapter{System requirements and specification}
In summary of the introduced concepts, we outline on a high level how an inspired system, and evaluation could arise from the above findings:

\paragraph{Evaluation} We would utilize the classification metric of the 2023 Google unnlearning competition, predicting images to age. Scores obtained would be compared to other submissions in the competition.

\paragraph{Cause Index} The latent cause theory model triangle relationship between conditioned stimuli, unconditioned stimuli, latent cause and conditional response. Would be reframed as data, unlearning marker, cause, and output. In which a model layer would be learnt to entirely dependent on the hidden causes in each layer to determine layer output. 

\paragraph{index resolution and unlearning} While cause index would serve as cause creation to the tree of causes in the same level, pattern seperation and completion would act as depth extension of the cause tree, in which higher layers would allow greater resolution. During unlearning, the causes would grow to allow direct pin-point of data, then removed or updated, to achieve unlearning.



\section{requirements gathering and prioritization}
\subsection{2023 unlearning competition metric}
    \paragraph{Link} 

    \url{https://github.com/Google-deepmind/unlearning_evaluation}

    \paragraph{Data} CASIA-CeFA \cite{liu2021casia,liu2021cross} \\
    Example label: .txt file on each line: <id> <year of birth> <F/M> \\
    Example target: .jpg 300*300 (variable size)

    \paragraph{Pre-trained model} 
    a ResNet-18 image classifier trained for 30 epochs on a dataset of face images, predicting age group (10 classes) \\
    Input: 32*32 (resized images)

    \paragraph{prioritization}
    High. Since without the metric, we cannot evaluate or compare methods of choice.

\subsection{latent cause theory model}
    \paragraph{Link} \url{https://github.com/sjgershm/memory-modification} \cite{gershman2017computational}
    \paragraph{prioritization}
    High. Core model design
\subsection{Pattern seperation}
    \paragraph{prioritization}
    High. Core model design
\subsection{Pattern completion}
    \paragraph{prioritization}
    Low. Optimization
\subsection{comparison with other models}
Compare against exact, fine-tune, and other unlearning methods.
    \paragraph{prioritization}
    Low. since unlearning metric score already used in large competition with precedence, plus running the inefficient unlearning metric on all of those models is infesible.

% \section{appropriate software development life cycle methodologies for development projects}

\chapter{Legal, Social Ethical and Professional issues}
\section{Inherent bias of training data and the limits of unlearning}
Training data is inherently biased, and the absorption of these biases into an increasingly AI-generated content landscape poses a substantial challenge to keeping models fair and neutral. Machine unlearning has been proposed as a partial remedy: rather than retraining a model from scratch on a re-curated dataset, biased or unrepresentative subsets of training data can, in principle, be selectively removed from an already-deployed model. In practice this is difficult. Bias is rarely localized to a small, identifiable set of training examples; it is typically distributed across the dataset and entangled with the model's learned representations, meaning that unlearning a specific subset may not fully remove its influence on downstream predictions. There is also a risk of \textit{catastrophic forgetting}, where aggressive unlearning degrades unrelated capabilities or, paradoxically, introduces new biases by shifting the model's decision boundaries in unpredictable ways. This raises an open research question central to the chapter: can unlearning be used not only as a compliance mechanism for erasure requests, but as a genuine tool for bias correction, and if so, how is that correction verified?

\section{Data sovereignty, intellectual property, and the right to erasure}
Data ownership and data security are fundamental to personal well-being, safety, and an accountable justice system, and machine unlearning sits at the technical intersection of these concerns. Where data protection regulations grant individuals a right to erasure, and where copyright holders seek removal of unlicensed material absorbed into a model's parameters, unlearning is the mechanism by which such rights could be practically enforced against a trained model rather than only against a raw dataset. This distinction matters: deleting a data point from a training corpus does nothing to remove its influence from a model already trained on it, so data sovereignty in the age of machine learning cannot be satisfied by database-level deletion alone. The same technical gap applies to intellectual property claims, where a rights holder may demand not just that infringing material be excluded from future training runs, but that its influence be actively unlearned from existing deployed models. This creates a tension between individual/creator control over data and the practical, computational cost of enforcing that control at the model level.

\section{Accountability and model transparency in unlearning}
Models are increasingly capable of producing results far exceeding human performance. Without a baked-in, human-readable, and accountable step chain, AI could potentially override human input and automate decisions without oversight aligned to common human interests. This concern extends directly to machine unlearning itself: an unlearning operation is only as trustworthy as the ability to verify that it actually occurred. Unlike conventional deletion, where removal from storage is straightforward to confirm, unlearning claims are made about a model's internal parameters, which are not human-interpretable by default. This creates an accountability gap: an organisation can claim compliance with an erasure request while having no rigorous, auditable proof that the targeted data's influence has actually been removed. Verifiable and certifiable unlearning, mechanisms that allow independent auditors or regulators to confirm that specific data no longer influences a model's outputs, is therefore a prerequisite for genuine accountability, not an optional technical nicety.

\chapter{System Design}
\section{Adaptation of latent cause theory model}
\subsection{Motivation for the Mapping}
 
In the original formulation, an animal's conditioned response on a given
trial depends entirely on which latent cause it currently believes is
active: the response is computed as a cause-weighted combination of
learned stimulus--outcome associations,
\[
V_t = \left(X_t^\top W\right) \cdot \mathrm{post}_t,
\]
where $W$ holds the per-cause learned weights and $\mathrm{post}_t$ is the
posterior probability over causes at trial $t$. Nothing about this
computation is specific to fear conditioning: it is a general statement
that a layer's output is a function of its input and of a hidden,
categorical state that the layer infers online. The reframing in
Table~\ref{tab:mapping} makes this explicit by relabeling the unconditioned
stimulus $r$ as an \emph{unlearning marker} --- a signal, present on some
trials and absent on others, that drives the model to either reinforce or
actively erase (via the protein-synthesis-inhibitor mechanism, \texttt{psi})
its current cause-conditional associations --- and by relabeling the
conditioned response as a generic \emph{output} that depends only on
\texttt{data} and the inferred \texttt{cause}.
 
\subsection{Implementation: \texttt{LatentCauseLayer}}
 
The reframing was implemented as a thin wrapper class,
\texttt{LatentCauseLayer}, around the unmodified core inference routine.
Its public interface exposes exactly the four renamed quantities:
 
\begin{verbatim}
class LatentCauseLayer:
    def forward(self, data, unlearning_marker, Dist):
        results = imm_localmap(data, unlearning_marker, Dist, self.opts)
        self.data = data
        self.unlearning_marker = unlearning_marker
        self.cause_posterior = results.Zp
        self.output = results.V
        self.causes = results.cause_index
        return self.output
\end{verbatim}
 
No numerical behavior of the underlying model was altered by this wrapper;
it exists purely to present the existing computation under the requested
vocabulary, and to expose the array described in
Section~\ref{sec:causearray}.
 
\section{The Growing Cause-Index Array}
\label{sec:causearray}
 
The original MATLAB implementation represents cause assignment with a
fixed-size matrix \texttt{Z} of shape $T \times K$, where $K$ is a
pre-allocated upper bound on the number of causes (default 15). Whether and
when a \emph{new} cause was created is not recorded anywhere; it can only
be inferred after the fact by inspecting which columns of \texttt{Z} are
nonzero. To satisfy the requirement that new causes populate an explicit,
growing array, the core inference loop was instrumented as follows.
 
\subsection{Definition of a Creation Event}
 
At the end of processing trial $t$, the model selects the maximum a
posteriori cause,
\[
k = \arg\max_j \ \mathrm{post}_t[j],
\]
where $\mathrm{post}_t$ is the trial's cause posterior after incorporating
the unlearning marker (reward). A creation event is defined precisely as
the case where cluster $k$ had a prior count of zero before trial $t$,
i.e. $N[k] = 0$, where $N$ is the vector of per-cause trial counts
accumulated from $Z[0{:}t,:]$. This is the same condition the ddCRP prior
already uses internally to reserve probability mass for an as-yet-unused
cause, so no new inference logic was introduced --- only a record of when
that branch is actually taken.
 
\subsection{Implementation}
 
\begin{verbatim}
cause_index = []  # grows by one entry per creation event
 
...
    k = int(np.argmax(post))
    if N[k] == 0:
        cause_index.append({'trial': t, 'cause_id': k})
    Z[t, k] = 1
\end{verbatim}
 
The list \texttt{cause\_index} starts empty and accumulates one dictionary
per creation event, each recording the trial index at which the cause was
created and the cause's integer identifier. It is attached to the
inference results as \texttt{results.cause\_index} and surfaced on the
layer wrapper as \texttt{layer.causes}, with convenience accessors
\texttt{layer.num\_causes\_created} and \texttt{layer.cause\_timeline()}
for extracting the trial indices and cause identifiers as parallel arrays.
 
\subsection{Demonstration}
 
A minimal standalone example --- three rewarded trials followed, after a
20-unit time gap, by five unrewarded trials --- was used to confirm the
array's behavior. With the learning-rate and threshold parameters used
elsewhere in the codebase's own experiment designs
($\eta = 0.3$, $\theta = 0.02$, $\lambda = 0.01$), the model creates a
first cause at trial 0 and, once the unlearning marker fails to match what
the first cause predicts, a second cause at trial 3:
\begin{verbatim}
cause creation events = [{'trial': 0, 'cause_id': 0},
                         {'trial': 3, 'cause_id': 1}]
\end{verbatim}

This matches the qualitative mechanism described in
Section~\ref{sec:causearray}: the first cause's learned weight, trained
toward a reward of 1 during acquisition, produces a large reward-prediction
mismatch as soon as the unlearning marker drops to 0, which is precisely
the condition that favors assigning MAP credit to a fresh, as-yet-unused
cause rather than continuing to update the existing one.

\section{Adaptation of unlearning metric}
\subsection{data}
The data was obtained from \url{https://sites.google.com/view/face-anti-spoofing-challenge/dataset-download/casia-surf-cefacvpr2020}, consisting of significantly more data of different format, from the original dataset, including three distinct geographies, masking, lighting and other irrelevant data to the original metric, significantly increasing training time. Adaptation considerations are listed as follows:

\paragraph{Dataset dimensions} Resize is changed from 224*224 to 16*16
\paragraph{label loader} SURFDataset class \_load\_labels are modified to read all .txt files to match labels with the most inner .jpg files under /color directory
\paragraph{age computation} Age computation from year of birth
\paragraph{GPU/Device support} Added Metal Performance Shaders support and correct floating point metal assumptions.
\paragraph{Class Weights Refactoring} List based to explicit tensor allocation of classes to handle missing classes case.
\paragraph{ID} Changed int resolve problems on GPU to string
\paragraph{testing} translated xmanager workflow to bash ready script in ./run.sh and more logging.info use to track progress in training models


% \section{Phases of development and experiments}
% \section{add details of design or experiments - challenges, problems, workarounds, adaptations, extensions, etc.}
% \section{Justification of design choices and highlight the challenges}

% \chapter{Experimental Results / Evaluation}
% \section{evaluate how the system performs}
% \section{discuss results and show insights into why you got that result}

\chapter{Conclusion and Discussion}
\section{Overall project conclusion}
We have been able to infer properties of reconsolidation therapy into a theoretical model, with complications it's implementation. Most of it stem from the fact that a complicated model requires complicated verification, even if it produces a runtime friendly output model.

We have been able to expreiment with compoents of our theoretical model. Namely latent cause model to infer cause based on experience and proximity, and metric.

The project's scope makes it difficult to complete in limited time, especially on complicated topics in reasoning about the vast choice of models, abstraction with no concrete detail of implementation or precedence. 

\section{future work}
To let a model truly grow and prune on reasoning, models with dynamic size would bridge the current knowledge gap in how the brain works and efficient implementation.
Nested Learning \cite{behrouz2025nested} might be useful where the rules of learning are continually discovered.
Memory and Data considerations might prove beneficial to consult self-modifying code, where data and programme act as expressions of limited resources in optimization. In this case the data would be causes and the programme it's expressions.

\bibliographystyle{agsm}
\bibliography{references}
\end{document}





